\subsection{稳定子集。诱导律}

定义6. 集合 E 的子集 A 称为在 E 上的合成律 \(\top\) 下稳定，如果 A 中两个元素的合成属于 A。此时映射 \((x,y)\mapsto x\top y\) （从 \(\mathbf{A}\times \mathbf{A}\) 到 A ）称为由律 \(\top\) 在 A 上诱导的律。带有由 \(\top\) 诱导的律的集合 A 称为 E 的子广群。

换言之，为使 A 在律 \(\top\) 下稳定，其充要条件是 \(A \top A \subset A\) 。E 的稳定子集与相应的子广群通常被视为同一事物。

E 的稳定子集的任一族的交集是稳定的；特别地，存在 E 的包含给定子集 X 的最小稳定子集 A；称 A 由 X 生成，而 X 称为 A 的生成系或生成集。相应的子广群也称为由 X 生成的。

命题1. 设E和F是两个广群，且f是E到F的一个同态。

\begin{enumerate}
\def\labelenumi{(\roman{enumi})}
\item
  在 \(f\) 下， \(\mathbf{E}\) 的稳定子集的像是 \(\mathbf{F}\) 的稳定子集.
\item
  在 \(f\) 下， \(\mathbf{F}\) 的稳定子集的逆像是 \(\mathbf{E}\) 的稳定子集.
\item
  设 \(\mathbf{X}\) 为 \(\mathbf{E}\) 的一个子集。在 \(f\) 下， \(\mathbf{E}\) 中由 \(\mathbf{X}\) 生成的稳定子集的像是 \(\mathbf{F}\) 中由 \(f(\mathbf{X})\) 生成的稳定子集.
\item
  若 \(g\) 是第二个同态，从 \(\mathbf{E}\) 到 \(\mathbf{F}\) ，则诸元素 \(x\) （属于 \(\mathbf{E}\) ，且使得 \(f(x) = g(x)\) ）所成的集合是 \(\mathbf{E}\) 的一个稳定子集.
\end{enumerate}

